\section{Fazit}
Für Lebensdauermessungen wurde eine Fast-Slow-Koinzidenzschaltung schrittweise aufgebaut. Zuerst wurde der Slow-Kanal, dann der Fast-Kanal kalibriert. Zur Energiekalibrierung wurden die Spektren der Natrium- und Bariumquellen aufgenommen und mithilfe von Gauß-Funktionen an die beobachteten Peaks angepasst. Die so ermittelten Peakpositionen wurden bekannten Übergangsenergien zugeordnet, wodurch die Energiekalibrierung des Mehrkanal-Koinzidenzanalysators (MCA) ermöglicht wurde. 
(Die geringe Genauigkeit der Anpassung ist auf kleine statistische und systematische Unsicherheiten zurückzuführen, die die Kalibrierung jedoch nicht wesentlich beeinträchtigten.)

Diese Kalibrierung ermöglichte auch die Bestimmung der Betriebsenergiebereiche des Einkanalanalysators (SCA) und des Konstantfraktionssplitters (CFD). Die beobachteten Felder entsprachen den erwarteten Linien: der 511-keV-Annihilationslinie im Natriumspektrum und den 81-keV- und 356-keV-Linien im Bariumspektrum. Zusätzlich wurde im Natriumspektrum das Compton-Kontinuum beobachtet, und im Bariumspektrum wurden weitere Verschiebungen, wie beispielsweise die 31-keV-Röntgenlinie, festgestellt.

Aus der 511-keV-Linie der Natriumquelle wurde eine schnelle Kurve erstellt. Die Zeitkalibrierung erfolgte anhand des Maximums. Die Zeitauflösung betrug:

$\overline{t}_\text{FWHM} = (\pm )\,\text{ns}.$


Basierend auf Lebensdauermessungen mit der Bariumquelle wurde die Lebensdauer des angeregten Zustands zu \(^{133}_{55}\text{Cs}\,(5/2^+)\) bestimmt. $[\tau = ( \pm) ns] $
Die zugehörige Halbwertszeit beträgt:
$T_{1/2} = (\pm)$ ns, was gut mit dem in der Literatur angegebenen Wert von $(\pm)$ ns (bei einer ?$\sigma$-Umgebung) übereinstimmt.